\SourceLeaf{070}
\Article{8}
Solutio generalis, quam hactenus explicavimus, ei potissimum casui adaptata est, ubi \textit{una} incognita a quantitatibus observatis pendens determinanda est. Quoties vero plures incognitae ab iisdem observationibus pendentes valores maxime plausibiles exspectant, vel quoties adhuc incertum est, quasnam potissimum incognitas ex observationibus derivare oporteat, has alia ratione praeparare conveniet, cuius evolutionem iam aggredimur.

Considerabimus quantitates \(x,y,z\) etc. tamquam indeterminatas, statuemus
\[
\left.\begin{aligned}\relax
[aa]x+[ab]y+[ac]z+\text{etc.}&=\xi\\
[ab]x+[bb]y+[bc]z+\text{etc.}&=\eta\\
[ac]x+[bc]y+[cc]z+\text{etc.}&=\zeta
\end{aligned}\right\}\ (6)
\]
etc., supponemusque, per eliminationem hinc sequi
\[
\left.\begin{aligned}\relax
[\alpha\alpha]\xi+[\alpha\beta]\eta+[\alpha\gamma]\zeta+\text{etc.}&=x\\
[\beta\alpha]\xi+[\beta\beta]\eta+[\beta\gamma]\zeta+\text{etc.}&=y\\
[\gamma\alpha]\xi+[\gamma\beta]\eta+[\gamma\gamma]\zeta+\text{etc.}&=z
\end{aligned}\right\}\ (7)
\]
etc.

Ante omnia hic observare oportet, coëfficientes symmetrice positos necessario aequales fieri, puta
\[
\begin{aligned}\relax
[\beta\alpha]&=[\alpha\beta]\\
[\gamma\alpha]&=[\alpha\gamma]\\
[\gamma\beta]&=[\beta\gamma]\\
\text{etc.}&
\end{aligned}
\]
quod quidem e theoria generali eliminationis in aequationibus linearibus sponte sequitur, sed etiam infra, absque illa, directe demonstrabitur.

Habebimus itaque
\[
\left.\begin{aligned}
x^0&=-[\alpha\alpha].[al]-[\alpha\beta].[bl]-[\alpha\gamma].[cl]-\text{etc.}\\
y^0&=-[\alpha\beta].[al]-[\beta\beta].[bl]-[\beta\gamma].[cl]-\text{etc.}\\
z^0&=-[\alpha\gamma].[al]-[\beta\gamma].[bl]-[\gamma\gamma].[cl]-\text{etc.}
\end{aligned}\right\}\ (8)
\]
\begin{center}etc.\end{center}
unde, si statuimus
\SourceLeaf{071}
\[
\left.\begin{aligned}\relax
[\alpha\alpha]\mathfrak A+[\alpha\beta]\mathfrak B+[\alpha\gamma]\mathfrak C+\text{etc.}&=A\\
[\alpha\beta]\mathfrak A+[\beta\beta]\mathfrak B+[\beta\gamma]\mathfrak C+\text{etc.}&=B\\
[\alpha\gamma]\mathfrak A+[\beta\gamma]\mathfrak B+[\gamma\gamma]\mathfrak C+\text{etc.}&=C
\end{aligned}\right\}\ (9)
\]
etc., obtinemus
\[
K=k-A[al]-B[bl]-C[cl]-\text{etc.}
\]
vel si insuper statuimus
\[
\left.\begin{aligned}
aA+bB+cC+\text{etc.}&=p\varepsilon\\
a'A+b'B+c'C+\text{etc.}&=p'\varepsilon'\\
a''A+b''B+c''C+\text{etc.}&=p''\varepsilon''
\end{aligned}\right\}\ (10)
\]
etc., erit
\[
K=k-l\varepsilon-l'\varepsilon'-l''\varepsilon''-\text{etc.}\qquad(11)
\]
\Article{9}
Comparatio aequationum (7), (9) docet, quantitates auxiliares \(A,B,C\) etc. esse valores indeterminatarum \(x,y,z\) etc. respondentes valoribus indeterminatarum \(\xi,\eta,\zeta\) etc. his \(\xi=\mathfrak A,\ \eta=\mathfrak B,\ \zeta=\mathfrak C\) etc., unde patet haberi
\[
\left.\begin{aligned}\relax
[aa]A+[ab]B+[ac]C+\text{etc.}&=\mathfrak A\\
[ab]A+[bb]B+[bc]C+\text{etc.}&=\mathfrak B\\
[ac]A+[bc]B+[cc]C+\text{etc.}&=\mathfrak C
\end{aligned}\right\}\ (12)
\]
etc. Multiplicando itaque aequationes (10) resp. per \(\frac ap,\ \frac{a'}{p'},\ \frac{a''}{p''}\) etc.\(\cdot\) et addendo, obtinemus
\[
\left.\begin{aligned}
\mathfrak A&=a\varepsilon+a'\varepsilon'+a''\varepsilon''+\text{etc.}\\
\text{et prorsus simili modo}\quad&\\
\mathfrak B&=b\varepsilon+b'\varepsilon'+b''\varepsilon''+\text{etc.}\\
\mathfrak C&=c\varepsilon+c'\varepsilon'+c''\varepsilon''+\text{etc.}
\end{aligned}\right\}\ (13)
\]
etc. Iam quum \(\mathfrak A\) sit valor functionis \(X\), si pro \(v,v',v''\) etc. valores observati substituuntur, facile perspicietur, si his applicentur correctiones \(-\varepsilon,\ -\varepsilon',\ -\varepsilon''\) etc. resp., functionem \(X\) hinc adepturam esse valorem \(0\), et perinde functiones
